\begin{corollary}[A normalized fifth-degree cancellation]
\label{powers:cor:canonical-fifth}
Use the fixed-coordinate finite parts
\(Q_k=\FP\int_0^1\psi(x)^k\,\dd x\), and put
\[
 \Omega_4=2D_{1,3}+D_{2,2}+2D_{3,1},
 \qquad D_{p,q}=[s]\zeta(p+s,q).
\]
Here \([s]\) denotes the coefficient of \(s\) in the Laurent
expansion, including when a pole is present. Define
\begin{align}
 P_5(T)={}&T^5+5\gamma T^4
      +10\bigl(\gamma^2+\zeta(2)\bigr)T^3\notag\\
 &\quad+10\bigl(\gamma^3+3\gamma\zeta(2)+\zeta(3)\bigr)T^2.
 \label{powers:eq:canonical-fifth-polynomial}
\end{align}
Then
\begin{align}
 \FP\int_0^1P_5\bigl(\psi(x)\bigr)\,\dd x
 ={}&-120\chi+10\Omega_4+10\zeta(2)\zeta'(2)\notag\\
 &+\frac{260}{3}\zeta(4)-60\gamma\zeta(4)
      -20\gamma^3\zeta(2)-20\zeta(2)\zeta(3).
 \label{powers:eq:canonical-fifth-finite-part}
\end{align}
Among monic polynomials
\(T^5+a_4T^4+a_3T^3+a_2T^2\) with coefficients in
\(\mathbb Q[\gamma,\zeta(2),\zeta(3)]\), the displayed polynomial is
the unique one whose coefficients of \(\theta,\eta,\gamma_1\)
vanish in the stated formulas for \(Q_2,Q_3,Q_4,Q_5\), when these
symbols are treated as formal coordinates.
\end{corollary}

\begin{proof}
The finite-part operation in one fixed coordinate is linear. In
\(Q_5+a_4Q_4+a_3Q_3+a_2Q_2\), cancellation of the displayed
\(\theta,\eta,\gamma_1\) coefficients gives, successively,
\begin{align*}
 120\gamma-24a_4&=0,\\
 60\bigl(\zeta(2)-\gamma^2\bigr)+24\gamma a_4-6a_3&=0,\\
 -20\bigl(\gamma^3-3\gamma\zeta(2)+\zeta(3)\bigr)
 +12\bigl(\gamma^2-\zeta(2)\bigr)a_4-6\gamma a_3+2a_2&=0.
\end{align*}
This triangular system has exactly the coefficients in
\eqref{powers:eq:canonical-fifth-polynomial}. The
\(\zeta'(3)\) coefficient simultaneously becomes
\(-20\gamma+4a_4=0\), while the coefficients of \(\chi\),
\(\Omega_4\), and \(\zeta(2)\zeta'(2)\) remain respectively
\(-120,10,10\). Before using any even-zeta reduction, the remaining
scalar is
\[
 \frac{35}{3}\zeta(4)+90\gamma\zeta(4)+30\zeta(2)^2
 -20\gamma^3\zeta(2)-60\gamma\zeta(2)^2
 -20\zeta(2)\zeta(3).
\]
Substitute \(\zeta(2)^2=\tfrac52\zeta(4)\) to obtain
\eqref{powers:eq:canonical-fifth-finite-part}.
\end{proof}

The uniqueness assertion is deliberately relative to the displayed
coefficient presentation and to the normalization excluding constant
and linear terms. It does not assert arithmetic independence of
\(\chi,\Omega_4,\theta,\eta,\gamma_1\), or the impossibility of
rewriting \(\Omega_4\) in another family of constants. In particular,
``cancellation of \(\gamma_1\)'' means cancellation of its explicit
coefficient while \(\Omega_4\) is retained as a coordinate; a
different representation of \(\Omega_4\) may contain
\(\gamma_1\). Further arithmetic reduction of the mixed-depth term
remains a separate question.
